% Methods: level-binned water frequency (stage 2 of the Earth Engine pipeline; A6).
% Design: D-015, D-027, D-060, D-062, D-068, D-071, D-074, D-076; PREREG sections 2-4,
% Amendments 3, 7 and 11. Every number comes from \metric{...}.

\subsection*{Level-binned water frequency}
\emph{Bins and halves.} Each scene was placed by its interpolated level in a bin \metric{params.bin_width_m}~m wide, counted from \metric{params.bin_origin_m}~m on the DAHITI reference; bins of \metric{params.bin_width_sens_lo_m} and \metric{params.bin_width_sens_hi_m}~m are a robustness check. Within each epoch and bin, acquisitions ordered in time were assigned alternately to two halves, the unit of the observational noise floor (pre-registration Amendment~3).

\emph{Counts.} Scenes enter if they fall in the open-water months, have scene cloud cover below \metric{params.max_scene_cloud_pct}\%, overlap the active zone, and have a level and a threshold. For every epoch, bin and half, Earth Engine counted at each pixel of a \metric{params.stage2_res_m}~m equal-area grid the acquisitions in which the pixel was water, water or dry, and classified at all (emergent vegetation included), after mosaicking the scenes of each acquisition. A scene whose predicted wind setup in a sector exceeded the screening threshold (Wind-setup screening) was left out in that sector. The counts cover the \metric{params.zone_tile_km}~km cells that contain the active zone, extended seaward by one cell, as insurance for bins below the levels that set the zone's seaward bound, and clipped to the domain (\metric{study_area.stage2_cover_km2}~\kmsq{}). Landsat and Sentinel-2 were exported separately from one scene plan, so that bins and halves are shared, and their counts summed.

\emph{Frequency.} The water frequency of a set of groups is the summed water count over the summed water-or-dry count, defined where the latter is at least \metric{params.min_valid_obs}. Groups combine by summing counts, never by averaging frequencies: the halves give a bin's full sample, and E2 and E3 together the sample of the topography. A cell is observed water at a bin if its frequency is at least \metric{params.contour_level}.

\emph{Emergent vegetation.} The emerged band is the set of cells in the domain that are water in the \metric{params.highstand_label} near-maximum extent and not in the \metric{params.lowstand_label} near-minimum extent of the active zone. A band cell is vegetated if emergent vegetation makes up at least \metric{params.veg_cell_share} of its classified E3 acquisitions (with at least \metric{params.min_valid_obs} of them), and a sector is vegetated if at least \metric{params.veg_sector_pct}\% of its band cells are; E3 describes the band as it is now. The Volga and Ural delta sectors are deltaic. The labels were fixed before any hindcast metric (pre-registration section~2 and Amendment~11).

\emph{Computation.} The pipeline ran in Google Earth Engine\citep{gorelick2017} in three stages: the active zone; a per-scene table of thresholds, histograms and masked fractions; and the per-pixel counts with the stored thresholds. Exports were piloted on small areas and metered against a budget while they ran. Counts were computed at \metric{params.stage2_res_m}~m rather than at the native Landsat resolution because their cost falls steeply with resolution; native-resolution counts over a sample area are a robustness check.
